\documentclass[12pt]{article}
\usepackage[T1]{fontenc}
\usepackage[utf8]{inputenc}
\usepackage{lmodern}
\usepackage{textcomp}
\usepackage{enumitem}
\usepackage{geometry}
\geometry{margin=1in}
\begin{document}
\begin{center}
Answer Key - Religious Studies - K-12 Level Assessment 7 \
\small (Answers are ordered exactly as in the question paper.)
\end{center}

\section*{Multiple Choice}
\begin{enumerate}[label=\arabic*.]
\item \textbf{Answer:} B
\\ \textit{Options:} A) Madhva; B) Patanjali; C) Ramanuja; D) Aryabhatta
\\ \textit{Note:} The Yoga Sutras are attributed to Patanjali, who is recognized as the compiler and author of this foundational text on yoga philosophy and practice in ancient Indian tradition.
\item \textbf{Answer:} A
\\ \textit{Options:} A) Bushido; B) Samurai; C) Nembutsu; D) Gomadaki
\\ \textit{Note:} The correct answer is Bushido because it refers to the ethical code and way of life that the samurai, or "warriors" of feudal Japan, adhered to, emphasizing values such as honor, loyalty, and discipline, which were articulated by figures like Yamago Soko.
\item \textbf{Answer:} D
\\ \textit{Options:} A) Zazen; B) Tariki; C) Kami; D) Ikebana
\\ \textit{Note:} Ikebana is the traditional Japanese practice of flower arranging that emphasizes harmony, balance, and the natural beauty of flowers, making it the correct answer.
\item \textbf{Answer:} C
\\ \textit{Options:} A) Dao; B) Xunzi; C) Mencius; D) Confucius
\\ \textit{Note:} Mencius is considered a mystic because his teachings emphasize the concept of qi (or ch'i) as a vital life force that influences moral development and the innate goodness of human nature, distinguishing his philosophical approach within Confucianism.
\item \textbf{Answer:} B
\\ \textit{Options:} A) The Hijab; B) The Crescent; C) The Minaret; D) The Qur'an
\\ \textit{Note:} The crescent moon is widely recognized as a symbol of Islam due to its historical associations with Islamic empires and its prominent display on flags and architecture, making it a significant cultural emblem for many Muslims and viewed similarly by Europeans.
\end{enumerate}

\section*{True / False}
\begin{enumerate}[label=\arabic*.]
\setcounter{enumi}{5}
\item \textbf{Answer:} False
\\ \textit{Options:} A) True; B) False
\\ \textit{Note:} The title Dalai Lama actually means 'Ocean of Wisdom,' not 'Ocean of Love,' in Tibetan culture.
\item \textbf{Answer:} True
\\ \textit{Options:} A) True; B) False
\\ \textit{Note:} The Upanishads are considered philosophical texts because they delve into profound inquiries regarding the nature of existence, the self, and the ultimate reality (Brahman), reflecting the core philosophical themes of Indian thought.
\item \textbf{Answer:} True
\\ \textit{Options:} A) True; B) False
\\ \textit{Note:} During times of drought and hardship, ordinary people in the Han Dynasty sought the intervention of deities, including the Queen Mother of the West, as she was associated with fertility and abundance in Chinese mythology.
\item \textbf{Answer:} True
\\ \textit{Options:} A) True; B) False
\\ \textit{Note:} The New Testament is a collection of Christian texts that comprises a total of 27 individual books, including the Gospels, letters, and other writings attributed to early followers of Jesus.
\item \textbf{Answer:} False
\\ \textit{Options:} A) True; B) False
\\ \textit{Note:} The statement is false because, after the Bar Kochba revolt, the primary centers for Jewish development were mainly in Babylonia and not in Europe, as significant Jewish life and scholarship flourished primarily in the Babylonian region during that period.
\end{enumerate}

\section*{Long Form Response}
\begin{enumerate}[label=\arabic*.]
\setcounter{enumi}{10}
\item \textbf{Answer:} The Meiji government played a crucial role in establishing a national cult centered around the emperor, promoting the idea that the emperor was a divine figure connected to kami, or spirits. This was part of a broader effort to unify the nation and foster a sense of national identity following Japan's modernization. By linking the emperor with kami, the government elevated his status, encouraging loyalty and reverence among the people. This connection influenced Japanese society by intertwining Shinto beliefs with state ideology, leading to the integration of religious practices into national life. As a result, the emperor became a symbol of unity and strength, significantly impacting both religion and culture in Japan during the Meiji period and beyond.
\\ \textit{Note:} correct because it accurately describes how the Meiji government promoted the emperor as a divine figure linked to kami to unify the nation and instill a sense of national identity, thereby influencing societal values and religious practices in Japan.
\item \textbf{Answer:} The goddess known as the Holy Mother of Mount Fairy Peach holds significant importance in the culture of Silla, one of the Three Kingdoms of Korea. Her worship reflects the values and beliefs of Silla society, particularly their deep connection to nature and the reverence for feminine divine figures. The Holy Mother symbolizes fertility, protection, and the nurturing aspects of the earth, which were vital for agricultural communities. Additionally, her veneration highlights the Silla people's belief in the balance between human life and the natural world, emphasizing harmony and respect for both. Overall, the worship of the Holy Mother of Mount Fairy Peach illustrates how the Silla people integrated their spiritual beliefs with their daily lives and environmental practices.
\\ \textit{Note:} correct because it identifies the Holy Mother of Mount Fairy Peach as a significant goddess in Silla culture, emphasizing her association with fertility, protection, and the nurturing aspects of nature, which reflect the agricultural values and feminine reverence central to Silla society.
\end{enumerate}

\vfill
\noindent\textit{End of Answer Key}
\end{document}